\section{Билет 44. Первое начало термодинамики. Элементарная и полная работа. Графический смысл работы газа. Работа идеального газа при изотермическом процессе (с выводом).}

\begin{definition}
В соответствии с \textbf{первым началом термодинамики}, теплота $Q$, полученная системой, расходуется на изменение внутренней энергии $\Delta U$ и совершение работы $A$:
\begin{equation}
    Q = \Delta U + A.
\end{equation}
Элементарная работа газа: $\delta A = p \, dV$. Полная работа в процессе $1 \to 2$: $A = \int_{V_1}^{V_2} p(V) \, dV$. Графически работа равна площади под кривой процесса в координатах $(p, V)$.
\end{definition}

\begin{theorem}[Вывод работы в изотермическом процессе]
Изотермический процесс протекает при постоянной температуре: $T = \text{const}$. Для идеального газа согласно закону Бойля--Мариотта:
\begin{equation}
    p V = \text{const} = \nu R T \quad \Rightarrow \quad p(V) = \frac{\nu R T}{V}. \label{eq:isothermal_p_44}
\end{equation}
Подставляем выражение для давления \eqref{eq:isothermal_p_44} в формулу полной работы:
\begin{equation}
    A = \int_{V_1}^{V_2} \frac{\nu R T}{V} \, dV.
\end{equation}
Так как $\nu$, $R$ и $T$ постоянны, выносим их за знак интеграла:
\begin{equation}
    A = \nu R T \int_{V_1}^{V_2} \frac{dV}{V} = \nu R T \left[ \ln V \right]_{V_1}^{V_2}.
\end{equation}
Вычисляя определённый интеграл, получаем формулу для работы:
\begin{equation}
    A = \nu R T \ln \frac{V_2}{V_1}. \label{eq:work_isothermal_V}
\end{equation}
Используя соотношение $V_2/V_1 = p_1/p_2$ (из $p_1 V_1 = p_2 V_2$), работу можно выразить через давления:
\begin{equation}
    A = \nu R T \ln \frac{p_1}{p_2}. \label{eq:work_isothermal_p}
\end{equation}
\end{theorem}

\begin{property}[Первое начало для изотермического процесса]
Внутренняя энергия идеального газа зависит только от температуры: $U = \frac{i}{2} \nu R T$. Поскольку $T = \text{const}$, изменение внутренней энергии равно нулю:
\begin{equation}
    \Delta U = 0.
\end{equation}
Подставляя это в первое начало термодинамики, получаем:
\begin{equation}
    Q = A. \label{eq:first_law_iso}
\end{equation}
Вся теплота, сообщаемая газу в изотермическом процессе, полностью превращается в механическую работу. При изотермическом сжатии ($V_2 < V_1$) работа внешних сил полностью переходит в теплоту, отдаваемую окружающей среде ($Q < 0$).
\end{property}

\begin{remark}
\textbf{Графический смысл и особенности:}
\begin{itemize}
    \item На $(p, V)$-диаграмме изотерма является гиперболой $p \sim 1/V$. Работа равна площади под этой гиперболой между $V_1$ и $V_2$.
    \item При расширении от одного объёма до другого изотермическая работа больше адиабатической, так как при адиабате температура падает и давление убывает быстрее (кривая адиабаты круче изотермы).
    \item Формулы \eqref{eq:work_isothermal_V} и \eqref{eq:work_isothermal_p} широко используются в расчётах тепловых машин, компрессоров и пневматических систем.
\end{itemize}
\end{remark}

\begin{figure}[htbp]
    \centering
    \begin{tikzpicture}[
        scale=1.2,
        >=Stealth,
        font=\small,
        thick
    ]

    \node[font=\bfseries] at (3, 4.5) {Изотермический процесс ($T = \text{const}$)};

    % === ОСИ ===
    \draw[->] (0, 0) -- (5.5, 0) node[right, font=\normalsize] {$V$};
    \draw[->] (0, 0) -- (0, 4) node[above, font=\normalsize] {$p$};
    \node[below left] at (0, 0) {$O$};

    % === ИЗОТЕРМА (гипербола p = 5/V) ===
    % Ограниченный домен: от V=1.2 (p≈4.2) до V=5.0 (p=1.0)
    % Кривая аккуратно вписана в область графика
    \draw[blue, very thick, domain=1.2:5.0, samples=100]
        plot (\x, {5/\x});

    % === ТОЧКИ 1 И 2 ===
    \def\Vone{1.5}
    \def\Vtwo{4.0}
    \def\Pone{3.33}   % 5/1.5
    \def\Ptwo{1.25}   % 5/4.0

    \fill[red] (\Vone, \Pone) circle (3pt);
    \node[above left] at (\Vone, \Pone) {$1$};
    \fill[red] (\Vtwo, \Ptwo) circle (3pt);
    \node[below right] at (\Vtwo, \Ptwo) {$2$};

    % === ЗАШТРИХОВАННАЯ ПЛОЩАДЬ (работа) ===
    \fill[blue!20, opacity=0.5, domain=\Vone:\Vtwo, samples=50]
        plot (\x, {5/\x}) -- (\Vtwo, 0) -- (\Vone, 0) -- cycle;

    % === ВЕРТИКАЛЬНЫЕ ЛИНИИ К ОСИ V ===
    \draw[dashed, gray] (\Vone, 0) -- (\Vone, \Pone);
    \draw[dashed, gray] (\Vtwo, 0) -- (\Vtwo, \Ptwo);

    % === ПОДПИСИ V1 И V2 ===
    \node[below] at (\Vone, 0) {$V_1$};
    \node[below] at (\Vtwo, 0) {$V_2$};

    % === ПОДПИСИ p1 И p2 ===
    \node[left] at (0, \Pone) {$p_1$};
    \node[left] at (0, \Ptwo) {$p_2$};

    % === СТРЕЛКА НАПРАВЛЕНИЯ ПРОЦЕССА ===
    \draw[->, blue, thick] (2.5, 2.0) -- (3.0, 1.67);

    % === БЛОК СВОЙСТВ (справа сверху, не перекрывает заголовок) ===
    \node[draw, fill=green!10, rounded corners, align=left, font=\small, inner sep=4pt] at (4.2, 3.2)
        {Изотерма:\\$pV = \text{const}$\\[4pt]
         $\Delta U = 0$\\[4pt]
         $Q = A$};

    % === ФОРМУЛА РАБОТЫ (внизу по центру) ===
    \node[draw, fill=yellow!10, rounded corners, align=center, font=\small, inner sep=5pt] at (2.75, -1.0)
        {Площадь под гиперболой:\\[6pt]
         $A = \nu RT \ln\dfrac{V_2}{V_1} = \nu RT \ln\dfrac{p_1}{p_2}$};

    \end{tikzpicture}
    \caption{Изотермический процесс на pV-диаграмме. Кривая — гипербола $p = \nu RT/V$. Работа газа равна площади под гиперболой между $V_1$ и $V_2$ (заштрихованная область).}
    \label{fig:isothermal_process_fixed}
\end{figure}

